% !TeX root = ../概率论与数理统计笔记.tex
% 由主文档通过 \input 载入；请编译主文档。
\setcounter{chapter}{14} % 保留原书章号。
\chapter{拟合优度、列联表与非参数方法}
\label{chap:ps-nonparametric}

本章研究分布形状、类别关联及较少依赖具体分布形式的推断。
先修内容为统计检验与经验分布；学习目标是按数据类型和研究问题选择方法，并识别各方法实际检验的假设。

\section{卡方拟合优度检验}
\label{sec:ps-goodness-fit}
引入\emph{拟合优度检验}（goodness-of-fit test），由类别观测频数与假设期望频数构造皮尔逊卡方统计量（Pearson chi-square statistic）。
区分完全给定的类别概率与由同一数据估计参数的模型，说明自由度调整所需的正则条件。
讨论期望频数较小时近似可能失效、类别合并须有依据，以及精确或模拟校准的替代方法。

\section{列联表中的独立性与同质性}
\label{sec:ps-contingency}
介绍\emph{列联表}（contingency table），比较独立性检验与同质性检验的抽样机制和原假设。
推导期望频数、约束与自由度；对小样本二维列联表介绍费舍尔精确检验（Fisher's exact test）的条件化思想。
结合条件概率说明统计关联不能单独确立因果关系。

\section{基于经验分布的检验}
\label{sec:ps-empirical-tests}
介绍柯尔莫哥洛夫--斯米尔诺夫检验（Kolmogorov--Smirnov test）比较经验分布与理论分布的思路。
标准单样本分布无关校准以完全给定的连续原假设分布为前提；由样本估计参数或出现离散取值时须重新校准。
说明图形比较、检验统计量和显著性结论各自承担的任务。

\section{符号检验、秩检验与置换思想}
\label{sec:ps-rank-tests}
介绍\emph{非参数方法}（nonparametric methods）中的符号检验、威尔科克森符号秩检验（Wilcoxon signed-rank test）及秩和检验（rank-sum test）。
分别说明中位数、对称位置模型与两样本分布比较的假设，不能把秩和检验无条件解释为中位数相等检验。
引入\emph{置换检验}（permutation test），以原假设下的可交换性或随机分组机制说明置换的合法性，并处理并列值。

\section{例题、习题与学习检查}
规划多类别比例是否符合指定模型、类别变量是否独立、配对差的符号检验以及两独立样本的秩比较。
习题要求根据抽样设计说明检验对象和校准方法，并比较参数方法与非参数方法的条件。
学习检查强调“非参数”仍包含抽样、独立性、对称性或可交换性等假设。





